% Part: incompleteness
% Chapter: representability-in-q
% Section: minimization-representable

\documentclass[../../../include/open-logic-section]{subfiles}

\begin{document}

\olfileid{inc}{req}{min}
\olsection{正则极小化在 $\Th{Q}$ 中是可表示的}

我们来考虑无界搜索。假设 $g(x, z)$ 是正则的并且在 $\Th{Q}$ 中可表示，例如由公式~$!A_g(x, z, y)$ 表示。令 $f$ 由 $f(z) = \umin{x}{[g(x, z) = 0]}$ 定义。我们想要找到一个表示~$f$ 的公式~$!A_f(z, y)$。$f(z)$ 的值是满足 (a) $g(x, z) = 0$ 且为其中最小的数~$x$，即对任意 $w < x$，$g(w, z) \neq 0$。因此，下面是一个自然的选择：
\[
!A_f(z,y) \ident !A_g(y, z, \Obj 0) \land \lforall[w][(w < y \lif \lnot
  !A_g(w, z, \Obj 0))].
\]
当然，在一般情形下，我们需要将 $z$ 替换为 $z_0$，……，$z_k$。

证明同样会涉及一些引理，表明 $\Th{Q}$ 足够强以致能证明某些事实。

\begin{lem}
\ollabel{lem:succ} 对每个常项符号~$a$ 和每个自然数~$n$，
\[
\Th{Q} \Proves \eq[(a' + \num n)][(a + \num n)'].
\]
\end{lem}

\begin{proof}
照例，证明通过对 $n$ 归纳进行。在基础情况 $n = 0$，我们需要证明 $\Th{Q}$ 可证明 $\eq[(a' + \Obj 0)][(a + \Obj
0)']$。但我们有：
\begin{align}
  \Th{Q} & \Proves \eq[(a' + \Obj 0)][a'] \quad \text{由公理 $Q_4$}
  \ollabel{step1}\\
  \Th{Q} & \Proves  \eq[(a + \Obj 0)][a] \quad \text{由公理 $Q_4$}
  \ollabel{step2} \\
  \Th{Q} & \Proves \eq[(a + \Obj 0)'][a'] \quad \text{由 \olref{step2}}
  \ollabel{step3} \\
  \Th{Q} & \Proves \eq[(a' + \Obj 0)][(a + \Obj 0)'] \quad
  \text{由 \olref{step1} 与 \olref{step3}}\notag
\end{align}
在归纳步中，可设已证 $\Th{Q}
\Proves \eq[(a' + \num n)][(a + \num n)']$。由于 $\num{n+1}$ 就是 $\num{n}'$，我们需要证明 $\Th{Q}$ 可证明 $\eq[(a' + \num{n}')][(a + \num{n}')']$。我们有：
\begin{align}
  \Th{Q} & \Proves \eq[(a' + \num n')][(a' + \num n)'] \quad
  \text{由公理 $!Q_5$} \ollabel{step5}\\
  \Th{Q} & \Proves \eq[(a' + \num n')][(a + \num n')'] \quad
  \text{归纳假设} \ollabel{step6}\\
  \Th{Q} & \Proves \eq[(a' + \num n)'][(a + \num n')'] \quad
  \text{由 \olref{step5} 和 \olref{step6}} \notag
\end{align}
\end{proof}

再次值得一提的是，这比说 $\Th{Q}$ 可证明 $\lforall[x][\lforall[y][(x' + y) = (x + y)']]$ 要弱。
尽管这个语句在~$\Struct{N}$ 中为真，但 $\Th{Q}$ 不能证明它。

\begin{lem}
\ollabel{lem:less-zero}
$\Th{Q} \Proves \lforall[x][\lnot x < \Obj 0]$。
\end{lem}

\begin{proof}
我们非形式化地给出证明（即，仅给出如何构造形式推导的提示）。

我们要对任意的~$a$ 证明 $\lnot a < \Obj 0$。根据 $<$ 的定义，我们需要在 $\Th{Q}$ 中证明 $\lnot
\lexists[y][\eq[(y'+a)][\Obj 0]]$。我们将假设 $\lexists[y][\eq[(y'+a)][\Obj 0]]$ 并导出一个矛盾。假设 $\eq[(b'+a)][\Obj 0]$。利用 $!Q_3$，我们可知 $\eq[a][\Obj 0] \lor \lexists[y][\eq[a][y']]$。我们区分情况。 

情况 1： $\eq[a][\Obj 0]$。由 $\eq[(b'+a)][\Obj 0]$，我们有 $\eq[(b' + \Obj 0)][\Obj 0]$。根据 $\Th{Q}$ 的公理~$!Q_4$，我们有 $\eq[(b' + \Obj 0)][b']$，因此 $\eq[b'][\Obj 0]$。但根据公理~$!Q_2$，我们也有 $\eq/[b'][\Obj 0]$，矛盾。

情况 2： 对某个 $c$，$\eq[a][c']$。那么我们有 $\eq[(b' +
c')][\Obj 0]$。由公理~$!Q_5$，我们有 $\eq[(b' + c)'][\Obj 0]$，这再次与公理~$!Q_2$ 矛盾。
\end{proof}

\begin{lem}
\ollabel{lem:less-nsucc}
对每个自然数~$n$，
  \[
  \Th{Q} \Proves
  \lforall[x][(x < \num {n+1} \lif (\eq[x][\Obj 0] \lor \dots \lor
    \eq[x][\num n]))].
  \]
\end{lem}

\begin{proof}
我们对~$n$ 使用归纳法。先考虑基础情况 $n = 0$。此时，我们需要对任意的~$a$ 证明 $a < \num 1 \lif \eq[a][\Obj 0]$。假设 $a < \num 1$。那么根据 $<$ 的定义公理，我们有 $\lexists[y][\eq[(y'+a)][\Obj 0']]$（因为 $\num 1 \ident \Obj 0'$）。

假设 $b$ 具有此性质，即 $\eq[(b'+a)][\Obj 0']$。我们需要证明 $\eq[a][\Obj 0]$。由公理~$!Q_3$，要么 $\eq[a][\Obj 0]$，要么存在 $c$ 使得 $\eq[a][c']$。前一种情况无需再证。故设 $\eq[a][c']$。那么我们有 $\eq[(b' + c')][\Obj 0']$。根据 $\Th{Q}$ 的公理~$!Q_5$，我们有 $\eq[(b'+c)'][\Obj 0']$。由公理~$!Q_1$，我们有 $\eq[(b' + c)][\Obj
0]$。但这意味着，根据公理~$!Q_8$，有 $c < \Obj 0$，与 \olref{lem:less-zero} 矛盾。

现在考虑归纳步。假设对~$n$ 成立，证明对 $n+1$ 的情况。假设 $a < \num {n+2}$。再次利用 $!Q_3$，我们可以区分两种情况：$\eq[a][\Obj 0]$，以及对某个 $b$，$\eq[a][b']$。在第一种情况，$\eq[a][\Obj 0] \lor \dots \lor
\eq[a][\num{n+1}]$ 平凡成立。在第二种情况，我们有 $b' < \num {n+2}$，即 $b' < \num{n+1}'$。由公理~$!Q_8$，对某个 $c$，$\eq[(c'+b')][\num{n+1}']$。由公理 $!Q_5$，$\eq[(c'+b)'][\num{n+1}']$。由公理~$!Q_1$，$\eq[(c'+b)][\num{n+1}]$，所以根据公理~$!Q_8$ 有 $b < \num{n+1}$。由归纳假设，$\eq[b][\Obj 0] \lor \dots \lor \eq[b][\num{n}]$。由此，由逻辑可得 $\eq[b'][\Obj 0'] \lor \dots \lor \eq[b'][\num{n}']$，又因 $\eq[a][b']$，所以得到 $\eq[a][\num{1}] \lor \dots \lor \eq[a][\num{n+1}]$。
\end{proof}

\begin{lem}
  \ollabel{lem:trichotomy} 对每个自然数~$m$，
  \[
  \Th{Q} \Proves
  \lforall[y][((y < \num{m} \lor \num{m} < y) \lor \eq[y][\num{m}])].
  \]
\end{lem}

\begin{proof}
对~$m$ 进行归纳。首先考虑 $m=0$ 的情形。由~$!Q_3$ 可得 $\Th{Q} \Proves
\lforall[y][(\eq[y][\Obj 0] \lor \lexists[z][\eq[y][z']])]$。
设 $a$ 任意。则要么 $\eq[a][\Obj 0]$，要么对某个~$b$，$\eq[a][b']$。在前一种情况下，我们也有 $(a < \Obj 0 \lor \Obj
0 < a) \lor \eq[a][\Obj 0]$。但若 $\eq[a][b']$，则由等词的逻辑可得 $\eq[(b' +
\Obj 0)][(a + \Obj 0)]$。由 $!Q_4$，$\eq[(a +
\Obj 0)][a]$，所以我们有 $\eq[(b' + \Obj 0)][a]$，从而得到 $\lexists[z][\eq[(z' + \Obj 0)][a]]$。根据 $!Q_8$ 中对 $<$ 的定义，有 $\Obj 0 < a$。如果 $\Obj 0 < a$，那么也有 $(\Obj 0 < a \lor a
< \Obj 0) \lor \eq[a][\Obj 0]$。 

现在假设我们有
\begin{align*}
  \Th{Q} & \Proves \lforall[y][((y < \num{m} \lor \num{m} < y) \lor
    \eq[y][\num{m}])]
  \intertext{而我们要证}
  \Th{Q} & \Proves \lforall[y][((y < \num{m+1} \lor \num{m+1} < y) \lor
    \eq[y][\num{m+1}])]
\end{align*}
设 $a$ 任意。由 $!Q_3$，要么 $\eq[a][\Obj 0]$，要么对某个~$b$，$\eq[a][b']$。在第一种情况下，由 $!Q_4$ 可得 $\eq[\num{m}' +
a][\num{m+1}]$，于是根据 $!Q_8$ 有 $a < \num{m+1}$。

现在考虑第二种情况，$\eq[a][b']$。由归纳假设，$(b < \num{m} \lor \num{m} < b) \lor
    \eq[b][\num{m}]$。

第一个析取支 $b < \num{m}$（由 $!Q_8$）等价于 $\lexists[z][\eq[(z' + b)][\num{m}]]$。假设 $c$ 具有此性质。
如果 $\eq[(c' + b)][\num{m}]$，那么也有 $\eq[(c' + b)'][\num{m}']$。由 $!Q_5$，$\eq[(c' + b)'][(c' + b')]$。因此，$\eq[(c' +
b')][\num{m}']$。对~$c'$ 进行存在概括，并注意到 $\num{m}'
\ident \num{m+1}$，可得 $\lexists[u][\eq[(u' + b')][\num{m+1}]]$。因此，如果 $b < \num{m}$，那么 $b' < \num{m+1}$，从而 $a < \num{m+1}$。

现在假设 $\num{m} < b$，即 $\lexists[z][\eq[(z' +
\num{m})][b]]$。假设 $c$ 是这样的 $z$，即 $\eq[(c' +
\num{m})][b]$。由逻辑，$\eq[(c' + \num{m})'][b']$。由 $!Q_5$，$\eq[(c' + \num{m}')][b']$。由于 $\eq[a][b']$ 且 $\num{m}' \ident
\num{m+1}$，$\eq[(c' + \num{m+1})][a]$。由 $!Q_8$，$\num{m+1} < a$。

最后，假设 $\eq[b][\num{m}]$。那么，由逻辑，$\eq[b'][\num{m}']$，于是 $\eq[a][\num{m+1}]$。

因此，从 $m$ 和 $b$ 的情形的每个析取支，我们都能得到 $m+1$ 和 $a$ 的对应析取支。
\end{proof}
  
\begin{prop}
\ollabel{prop:rep-minimization}
如果 $!A_g(x, z, y)$ 在~$\Th{Q}$ 中表示 $g(x, z)$，那么
\[
!A_f(z,y) \ident !A_g(y, z, \Obj 0) \land \lforall[w][(w < y \lif \lnot
  !A_g(w, z, \Obj 0))]
\]
表示 $f(z) = \umin{x}{[g(x, z) = 0]}$。
\end{prop}

\begin{proof}
首先证明，如果 $f(n) = m$，那么 $\Th{Q} \Proves !A_f(\num n, \num m)$，即  
\begin{align}
  \Th{Q} & \Proves !A_g(\num{m}, \num{n}, \Obj 0) \land \lforall[w][(w <
    \num{m} \lif \lnot !A_g(w, \num{n}, \Obj 0))]. \notag
  \intertext{由于 $!A_g(x, z, y)$ 表示 $g(x, z)$ 且当 $f(n) = m$ 时 $g(m, n) = 0$，我们有}
\Th{Q} & \Proves !A_g(\num{m}, \num{n}, \Obj 0). \notag
\intertext{若 $f(n) = m$，则对每个 $k < m$，$g(k, n) \neq 0$。所以}
\Th{Q} & \Proves \lnot !A_g(\num{k}, \num{n}, \Obj 0). \notag
\intertext{我们得到}
\Th{Q} & \Proves \lforall[w][(w < \num{m} \lif \lnot
  !A_g(w, \num{n}, \Obj 0))]. \ollabel{rep-less}
\end{align}
在 $m = 0$ 时由 \olref{lem:less-zero} 可得，否则由 \olref{lem:less-nsucc} 可得。

现在证明，如果 $f(n) = m$，那么 $\Th{Q} \Proves
\lforall[y][(!A_f(\num{n}, y) \lif \eq[y][\num{m}])]$。我们再次非形式化地概述该论证，将形式化工作留给读者。

假设 $!A_f(\num{n}, b)$。由此我们得到 $!A_g(b, \num{n},
\Obj 0)$ 和 $\lforall[w][(w < b \lif \lnot !A_g(w, \num{n}, \Obj
  0))]$。由 \olref{lem:trichotomy}，$(b < \num{m} \lor \num{m} < b)
\lor \eq[b][\num{m}]$。我们将证明 $b < \num{m}$ 和 $\num{m} < b$ 都会导致矛盾。

如果 $\num{m} < b$，那么由 可得 $\lnot !A_g(\num{m}, \num{n}, \Obj 0)$。但 $m = f(n)$，所以 $g(m, n) = 0$；又因 $!A_g$ 表示~$g$，所以 $\Th{Q} \Proves
!A_g(\num{m}, \num{n}, \Obj 0)$。于是得出矛盾。

现在假设 $b < \num{m}$。那么由 \olref{rep-less}，因 $\Th{Q} \Proves \lforall[w][(w <
  \num{m} \lif \lnot !A_g(w, \num{n}, \Obj 0))]$，可得 $\lnot !A_g(b, \num{n}, \Obj 0)$。这又与 矛盾。
\end{proof}

\end{document}